% 07_matrix.tex — 综合对比与阈值矩阵

\section{综合对比与可行性翻转分析}

\subsection{跨轨道综合成本指数（M5 模型）}

基于 M1--M5 五个独立 Python 建模脚本，整合发射成本（30\%）、散热面积（20\%）、光伏需求（15\%）、辐射威胁（25\%）和硬件寿命（10\%）五个维度的加权综合指数。指数值越低表示综合可行性越优，LEO 设为基准 1.00。

\begin{table}[H]
\centering
\caption{跨轨道综合成本指数排名（M5 模型输出）}
\label{tab:m5-ranking}
\resizebox{\textwidth}{!}{%
\begin{tabular}{@{}lcccccc@{}}
\toprule
\textbf{排名} & \textbf{轨道} & \textbf{综合指数} & \textbf{评级} & \textbf{发射} & \textbf{散热} & \textbf{辐射} \\
\midrule
\textbf{1} & \textbf{GEO}     & \textbf{0.83} & A & 1.53x & 1,040 & 7.7/100 \\
2 & 地月 L1             & 0.94 & A & 1.37x & 1,046 & 13.1/100 \\
3 & 地月 L4/L5          & 0.94 & A & 1.40x & 1,038 & 13.1/100 \\
4 & 地月 L2             & 0.95 & A & 1.40x & 1,046 & 13.1/100 \\
5 & 地日 L1             & 0.99 & A & 1.36x & 1,038 & 16.7/100 \\
6 & 地日 L2             & 0.99 & A & 1.36x & 1,038 & 16.7/100 \\
7 & LEO (基准)          & 1.00 & A & 1.00x & 482  & 10.0/100 \\
8 & MEO                 & 1.47 & B & 1.31x & 1,040 & 35.0/100 \\
9 & HEO                 & 2.52 & C & 1.60x & 1,015 & 50.0/100 \\
\bottomrule
\end{tabular}%
}
\end{table}

\textbf{核心发现}：6 个非 LEO 轨道的综合指数低于或等于 LEO。GEO 以 0.83 排名第一（比 LEO 优 17\%），因为其散热效率（1,040 vs 482~\wpm）和光伏面积需求（3,315 vs 8,844 m\textsuperscript{2}）的优势充分抵消了 $\sim$1.5x 的发射成本溢价。

\begin{figure}[H]
\centering
% figures/m5_ranking.tex — M5 综合成本指数条形图
% 数据来源：M5 orbit_cost_matrix.py

\begin{tikzpicture}
\begin{axis}[
    width=0.85\textwidth,
    height=9cm,
    ybar,
    bar width=12pt,
    enlarge x limits=0.15,
    ymin=0, ymax=3.0,
    ylabel={综合成本指数（越低越优）},
    xlabel={},
    ytick={0,0.5,1.0,1.5,2.0,2.5,3.0},
    symbolic x coords={GEO,地月L1,地月L4/L5,地月L2,地日L1,地日L2,LEO,MEO,HEO},
    xtick=data,
    xticklabel style={rotate=45, anchor=east, font=\footnotesize},
    nodes near coords,
    nodes near coords align={vertical},
    nodes near coords style={font=\footnotesize},
    legend style={at={(0.5,-0.22)}, anchor=north, font=\footnotesize, legend columns=5},
    grid=major,
    major grid style={dashed,gray!30},
]
% 以 LEO=1.00 为基准线
\draw[red, thick, dashed] (rel axis cs:0,0) -- (rel axis cs:1,0) node[above left, font=\small] {LEO 基准};

\addplot[fill=geocolor]       coordinates {(GEO,0.83)};
\addplot[fill=eml1l2color]    coordinates {(地月L1,0.94)};
\addplot[fill=l4l5color]      coordinates {(地月L4/L5,0.94)};
\addplot[fill=eml1l2color]    coordinates {(地月L2,0.95)};
\addplot[fill=l1l2color]      coordinates {(地日L1,0.99)};
\addplot[fill=l1l2color]      coordinates {(地日L2,0.99)};
\addplot[fill=leoorange]      coordinates {(LEO,1.00)};
\addplot[fill=meocolor]       coordinates {(MEO,1.47)};
\addplot[fill=heocolor]       coordinates {(HEO,2.52)};

\addlegendentry{GEO}
\addlegendentry{地月L1}
\addlegendentry{地月L4/L5}
\addlegendentry{地月L2}
\addlegendentry{地日L1}
\addlegendentry{地日L2}
\addlegendentry{LEO（基准）}
\addlegendentry{MEO}
\addlegendentry{HEO}
\end{axis}
\end{tikzpicture}

\caption{跨轨道综合成本指数排名（M5 模型）。指数越低越优，LEO 为基准 1.00。}
\label{fig:m5-ranking}
\end{figure}

\subsection{辐射威胁排序：TID vs SEE}

\begin{table}[H]
\centering
\caption{各轨道辐射威胁排序（M2 模型输出）}
\label{tab:m2-ranking}
\resizebox{\textwidth}{!}{%
\begin{tabular}{@{}lcccccc@{}}
\toprule
\textbf{轨道} & \textbf{TID中值} & \textbf{SEU (log10)} & \textbf{TIDvsLEO} & \textbf{寿命(yr)} & \textbf{威胁分} & \textbf{主威胁} \\
\midrule
HEO  & 125.0~\kradyr & 5.50 & 6.25x & 0.8  & 50.0 & \textbf{TID+SEE双杀} \\
MEO  & 10.0~\kradyr  & 5.00 & 0.50x & 10.0 & 35.0 & \textbf{SEE主导} \\
地日 L1/L2 & 3.5~\kradyr & 2.50 & 0.17x & 28.6 & 16.7 & \textbf{SEE主导} \\
地月 L4/L5 & 10.0~\kradyr & 1.35 & 0.50x & 10.0 & 13.1 & 混合偏SEE \\
GEO  & 6.5~\kradyr  & 0.74 & 0.33x & 15.4 & 7.7  & 混合型 \\
LEO  & 20.0~\kradyr & 0.00 & 1.00x & 5.0  & 10.0 & \textbf{TID主导} \\
\bottomrule
\end{tabular}%
}

\smallskip
\noindent 威胁分 = TID(40\%) + SEU(60\%)，TID 失效阈值 100~\kradyr。寿命基于 TID 中值 / 100~\kradyr 估算。
\end{table}

\subsection{TID 悖论深度分析}

"TID 悖论"是本调研最重要的发现：

\begin{observation}[TID悖论]
深空（地日 L1/L2）在适度屏蔽（$\sim$2.5 mm Al）后的年 TID 仅为 2--5~\kradyr，而穿越 SAA 的 LEO 轨道（SSO、高倾角）年 TID 为 10--30~\kradyr——深空 TID 反而更低。但深空 SEU 率比 LEO 高 100--1000 倍。
\end{observation}

物理原因：LEO 的 TID 主要由内范艾伦带捕获质子贡献（SAA 区域），深空无此来源。但深空无地磁场偏转 GCR 重离子，导致 SEU 率比 LEO 高出 2--3 个数量级 \cite{jahan_radiation_2026}。

\textbf{对算力硬件的工程判断}：在拉格朗日点部署算力的辐射挑战是 \textbf{SEE 挑战，不是 TID 挑战}。只要 SEE 被充分应对（TMR + 快速 scrubbing + SEL 保护），TID 甚至不是制约因素（屏蔽后 $<$5~\kradyr，10 年仅 50 krad，FDSOI 22nm 工艺在 400 krad 下 Vth 漂移仅 0.052 V——接近免疫 \cite{jahan_radiation_2026}）。

\begin{figure}[H]
\centering
% figures/tid_vs_seu.tex — TID vs SEU 散点图（TID 悖论可视化）
% 数据来源：M2 radiation_dose.py

\begin{tikzpicture}
\begin{axis}[
    width=0.85\textwidth,
    height=9cm,
    xlabel={TID 中值 (\kradyr)},
    ylabel={SEU 相对风险 (log10, LEO=0)},
    xmin=0, xmax=140,
    ymin=-0.5, ymax=6.0,
    scatter/classes={
        leo={mark=square*,leoorange},
        meo={mark=triangle*,meocolor},
        geo={mark=diamond*,geocolor},
        heo={mark=pentagon*,heocolor},
        l1l2={mark=*,l1l2color},
        l4l5={mark=star,l4l5color}
    },
    grid=major,
    major grid style={dashed,gray!30},
    nodes near coords,
    every node near coord/.append style={font=\footnotesize, anchor=west},
]

% LEO 基准
\addplot[scatter, only marks, scatter src=explicit symbolic]
    coordinates {
        (20.0, 0.0) [leo]
    };
\node[anchor=south west, font=\footnotesize] at (axis cs:20.0,0.0) {LEO (基准)};

% MEO - SEE主导
\addplot[scatter, only marks, scatter src=explicit symbolic]
    coordinates {
        (10.0, 5.00) [meo]
    };
\node[anchor=west, font=\footnotesize] at (axis cs:10.0,5.00) {MEO};

% GEO - 混合型
\addplot[scatter, only marks, scatter src=explicit symbolic]
    coordinates {
        (6.5, 0.74) [geo]
    };
\node[anchor=south west, font=\footnotesize] at (axis cs:6.5,0.74) {GEO};

% HEO - 双杀
\addplot[scatter, only marks, scatter src=explicit symbolic]
    coordinates {
        (125.0, 5.50) [heo]
    };
\node[anchor=south east, font=\footnotesize] at (axis cs:125.0,5.50) {HEO};

% 地日L1/L2 - SEE主导 (低TID)
\addplot[scatter, only marks, scatter src=explicit symbolic]
    coordinates {
        (3.5, 2.50) [l1l2]
    };
\node[anchor=south east, font=\footnotesize, text=l1l2color] at (axis cs:3.5,2.50) {地日L1/L2};

% 地月L4/L5
\addplot[scatter, only marks, scatter src=explicit symbolic]
    coordinates {
        (10.0, 1.35) [l4l5]
    };
\node[anchor=east, font=\footnotesize, text=l4l5color] at (axis cs:10.0,1.35) {地月L4/L5};

\end{axis}
\end{tikzpicture}

\caption{TID vs SEU 二维分布图。横轴为 TID 中值（\kradyr），纵轴为 SEU 相对 LEO 的对数风险。地日 L1/L2 位于左下角"TID 悖论区"——TID 极低但 SEE 风险高。HEO 位于右上角"双杀区"——TID 与 SEE 均极高。}
\label{fig:tid-vs-seu}
\end{figure}

\subsection{SEE 防护开销量化}

\begin{table}[H]
\centering
\caption{SEE 防护策略开销量化}
\label{tab:see-cost}
\small
\begin{tabular}{@{}lccc@{}}
\toprule
\textbf{防护策略} & \textbf{面积/功耗开销} & \textbf{适用范围} & \textbf{深空约束} \\
\midrule
TMR (全模块)     & $\sim$200\% & 全部逻辑路径 & 不可缩减 \\
TMR (选择性)     & $\sim$56\% LUT + 28\% 寄存器 & 关键路径 & 需验证覆盖率 \\
EDAC (SECDED)    & 12.5\% 内存 & SRAM/DRAM & MBU 可能突破 \\
Scrubbing        & $\sim$1--5\% 系统功耗 & FPGA 配置内存 & 间隔需缩至秒级 \\
SEL 保护        & 电路级 & 电源管理 & 必需 \\
\bottomrule
\end{tabular}
\end{table}

在 L2 环境中（SEU 率 100--1000x LEO），TMR + EDAC + scrubbing 的综合开销估计为面积 $\sim$200--300\%、功耗 $\sim$150--200\%，相较于无防护的 COTS 硬件 \cite{kchaou_tmr_2026}。

\subsection{"慢死 vs 猝死"判断}

\begin{observation}[SEE 比 TID 更难应对]
TID 可通过工艺选择一次性解决（FDSOI、GaN、SOI），而 SEE 防护（TMR + scrubbing）是\textbf{持续性的面积/功耗/复杂度惩罚}，无法通过单次工程决策消除。此外，SEE 率受太阳活动 10x+ 调制，SPE 不可预测——裕度设计困难。结论：SEE（猝死）比 TID（慢死）更难应对。
\end{observation}

\subsection{可行性翻转阈值矩阵}

表~\ref{tab:flip-threshold} 汇总各轨道实现"TCO 与地面可比"的临界条件。硬件成本降低（抗辐射溢价缩减）是所有轨道的共同瓶颈——硬件成本占 TCO 的 50--70\%。

\begin{table}[H]
\centering
\caption{各轨道可行性翻转阈值矩阵}
\label{tab:flip-threshold}
\small
\begin{tabular}{@{}lccccc@{}}
\toprule
\textbf{轨道} & \textbf{发射成本临界值} & \textbf{硬件溢价临界值} & \textbf{寿命临界值} & \textbf{翻转概率(2035)} \\
\midrule
LEO    & $<$\$100/kg & $<$3x 地面 & $>$7 年  & 中低 (25--40\%) \\
MEO    & $<$\$80/kg  & $<$2x 地面 & $>$10 年 & \textbf{极低 ($<$10\%)} \\
GEO    & $<$\$120/kg & $<$3x 地面 & $>$5 年  & 中低 (20--35\%) \\
HEO    & $<$\$60/kg  & $<$2x 地面 & $>$12 年 & \textbf{极低 ($<$5\%)}  \\
L1/L2  & $<$\$100/kg & $<$2.5x 地面 & $>$8 年  & 低 (10--20\%) \\
L4/L5  & $<$\$100/kg & $<$2.5x 地面 & $>$7 年  & 低 (15--25\%) \\
\bottomrule
\end{tabular}
\end{table}

\textbf{关键判断}：GEO 拥有最宽松的翻转条件——15 年寿命已获产业验证，RSGS 可进一步延寿。\textbf{GEO 是唯一可能率先实现经济可行的非 LEO 轨道}。MEO 和 HEO 的翻转条件极为苛刻，在 2035 年前几乎不可能。

\subsection{1 MW 散热与光伏面积需求}

\begin{table}[H]
\centering
\caption{1 MW 算力散热面积与光伏面积需求（M3/M4 模型）}
\label{tab:area-requirement}
\small
\begin{tabular}{@{}lcccc@{}}
\toprule
\textbf{轨道} & \textbf{350K散热 m\textsuperscript{2}} & \textbf{光伏面积 m\textsuperscript{2}} & \textbf{储能需求} & \textbf{日照占比} \\
\midrule
LEO       & 3,115 & 8,844 & 有 & 60\%   \\
MEO       & 1,439 & 3,528 & 无 & 95\%   \\
GEO       & 1,442 & 3,315 & 极少 & 98\%   \\
HEO       & 1,477 & 4,407 & 有 & 85\%   \\
地日 L1/L2 & 1,445 & 3,184 & \textbf{零} & \textbf{100\%} \\
地月 L4/L5 & 1,445 & 3,184 & \textbf{零} & \textbf{100\%} \\
\bottomrule
\end{tabular}
\end{table}

\textbf{零储能轨道}：地日 L1/L2 和地月 L4/L5 享受连续日照，无需储能电池——简化设计、降低质量。
